% language=us runpath=texruns:manuals/moretocome

\startcomponent moretocome-integration

\environment moretocome-style

\startchapter[title={Integration}]

When you put text on a curve there are (at least) two variants to
consider:

\startitemize
    \startitem
        You typeset the text and put glyphs along the shape, rotating the glyph
        following the curve. Rotation happens around the the center of the
        baseline of the glyph.
    \stopitem
    \startitem
        You typeset the text and transform the outline according to the position
        on the shape. So, when the shape is a circle the glyphs become wider at
        their top.
    \stopitem
\stopitemize

The first variant uses the original glyph path and therefore we can use the
regular font resource, while the second variant has to fetch the outline from the
font resource and work with the path. Here we discuss the second case.

\startbuffer[1]
\startMPcode
    path shape ; shape := reverse fullcircle
        xyscaled (50,40)
    ;
    draw lmt_followtext [
        text = "\bf Do you see the difference? ",
        path = shape,
    ] ysized 5cm withcolor "middleblue" ;
\stopMPcode
\stopbuffer

\startbuffer[2]
\startMPcode
    path shape ;
    shape := reverse fullcircle
        xyscaled (500,400)
    ;
    path p ; p := lmt_oscillated [
        text = "\bf Do you see the difference? ",
        path = shape,
    ] ;
    fill p withcolor "middleblue" ;
\stopMPcode
\stopbuffer

\startplacefigure[reference=fig:integration:0]
    \startcombination
        {\scale[width=.4tw]{\getbuffer[1]}} {first variant}
        {\scale[width=.4tw]{\getbuffer[2]}} {second variant}
    \stopcombination
\stopplacefigure

Because we gradually translate points we cannot simply take the existing points
that define the glyph. For instance, the topmost line of the \type {T} normally
has two points and just translating these two ignores the fact that in the middle
we move up. As demonstrated in \in {figure} [fig:integration:1] we have to add
points.

\startbuffer
\startMPcode
    path shape ;
    shape := reverse subpath(1,3) of fullcircle
        xyscaled (500,400)
    ;
    draw shape
        withpen pencircle scaled 2
        withcolor "middlegray"
    ;
    path p ; p := lmt_oscillated [
        text       = "\bf T T T T ",
        path       = shape,
        resolution = 5,
    ] ;
    fill       p withcolor "middleyellow" ;
    drawpoints p withcolor "middlered" ;
\stopMPcode
\stopbuffer

\startplacefigure[reference=fig:integration:1]
    \scale[width=tw]{\getbuffer}
\stopplacefigure

In \CONTEXT\ speak, this example is generated with:

\typebuffer[option=TEX]

This feature is a good example of the way we integrate \TEX, \LUA\ and \METAPOST\
so let's see what happens here. We start at the \TEX\ end, where we ask for some
\METAPOST\ graphic. The code between \typ {\startMPcode} and \typ {\stopMPcode}
is picked up at the \LUA\ end and passed to a \METAPOST\ instance. When processed
the result is picked up by \LUA, and the graphic operations become \PDF\ code.

When we are in \METAPOST, we need the outline of the text. So there \METAPOST\
picks up the text and (via \LUA) passes it to \TEX\ where it gets typeset and
ends up in a box. Processing the text uses the (set) font and font features are
applied, think of ligature building and kerning. The resulting node list contains
references to glyphs in a font as well as kerns (and maybe rules). Because we
want the outlines, we trigger a dedicated so called \quote {shipout} routine, one
that assembles a picture from those outlines. That picture is then passed back to
\METAPOST, where successive paths get positioned along the shape whereby points
are added and relocated. Some glyphs are made from multiple paths, for instance
an \type {O} has two paths and a \type {B} has three.

Getting the outline from the font means that we need to collect the so called
\CFF\ or \TTF\ glyph path descriptions and translate them to \METAPOST\
operations. Reading a font and parsing these path descriptions is again a task
for \LUA. We actually need these (original) path descriptions when we subset a
font for embedding and have to manipulate them when we are dealing with a
variable font. The export to a \METAPOST\ variant was written at the time we
wrote the \LUA\ code, if only because we love tracing such properties and
features.

It will be clear that we call \LUA\ from \TEX\ as well as \METAPOST\ various
times. In this case \typ {lmt_oscillated} is rather simple, most code involves
calculating the (additional) points from given points, the rest was just there.
This macro actually calls out to \LUA\ where we pick up the parameters. Only when
we are done typesetting and collecting glyph paths we return to \METAPOST. We use
scanners to get data and injectors to push data at the \TEX\ and \METAPOST\ end.
These mechanism are quite efficient. We basically extend both languages with new
features. We can summarize the process as:

\startlinecorrection
\TEX      \low{1} \rightarrow
\LUA      \low{1} \rightarrow
\METAPOST \low{1} \rightarrow
\LUA      \low{2} \rightarrow
\TEX      \low{2} \rightarrow
\LUA      \low{3} \rightarrow
\LUA      \low{4} \rightarrow
\METAPOST \low{2} \rightarrow
\LUA      \low{5} \rightarrow
\PDF      \low{1} \rightarrow
\TEX      \low{3}%
\stoplinecorrection

The various steps shown here are:

\starttabulate
\NC \TEX      \low{1} \NC ask for graphic with macro \NC \NR
\NC \LUA      \low{1} \NC prepare graphic data and pass to instance \NC \NR
\NC \METAPOST \low{1} \NC process graphic and call macro \NC \NR
\NC \LUA      \low{2} \NC pick up arguments \NC \NR
\NC \TEX      \low{2} \NC typeset text \NC \NR
\NC \LUA      \low{3} \NC pickup outlines \NC \NR
\NC \LUA      \low{4} \NC convert result to \METAPOST\ shape \NC \NR
\NC \METAPOST \low{2} \NC position transformed glyphs \NC \NR
\NC \LUA      \low{5} \NC pickup result from run and convert \NC \NR
\NC \PDF      \low{1} \NC graphic description \NC \NR
\NC \TEX      \low{3} \NC position and scale result in document \NC \NR
\stoptabulate

We understand that writing this solution from scratch without any code being
present takes some effort. However, as were already using \LUATEX\ and
\LUAMETATEX\ for a while, so in the end we only had to spend time on step \METAPOST
\low{2}. Mikael Sundqvist and I started our explorations with a more complex
approach, which actually exposed some complications so eventually we decided on a
slightly more demanding variant in terms of runtime, and the final code actually
has less lines. \in {Figure} [fig:integration:6] demonstrates that ligatures are
handled well.

\startbuffer
\startMPcode
    path shape ;
    shape := reverse subpath(1,3) of fullcircle
        xyscaled (500,400)
    ;
    path p ; p := lmt_oscillated [
        text       = "\bf effe flink fietsen",
        path       = shape,
        resolution = 20,
    ] ;
    fill p withcolor "middleyellow" ;
    draw p withcolor "middleblue" withpen pencircle scaled 1 ;
\stopMPcode
\stopbuffer

\startplacefigure[reference=fig:integration:6]
    \scale[width=tw]{\getbuffer}
\stopplacefigure

Of course we can combine this with other integrated features, for instance with
so called \quote {Perlin Noise} which generates patterns that look random but
where distortions are related (think for instance clouds):

\startbuffer
\definecharacterkerning
  [mine]
  [factor=0.05]

\startMPcode
    draw lmt_noise [
        bytemap     = 4,
        nx          = 1000,
        ny          = 1000,
        nz          = 1,
        iterations  = 2,
        frequency   = 0.1,
        amplitude   = 1,
        persistence = 0.5,
        lacunarity  = 2.0,
        minimum     = 0,
        maximum     = 255,
        method      = "angle",
    ] xysized (4cm,4cm) shifted (-2cm,-2cm) ;
    fill lmt_oscillated [
        text = "\setcharacterkerning[mine1]\bf\CONTEXT\ LMTX is FUN! "
        path = reverse fullcircle xyscaled (25mm,25mm) ,
    ] withtransparency (1,.75) withcolor white ;
\stopMPcode
\stopbuffer

\typebuffer[option=TEX]

Here we first create a background: again we parse the specification from \LUA\
using \METAPOST\ scanning routines, then we generate a byte map with noise using
a small \CCODE\ library that we access from \LUA. When done we tell \METAPOST\
the dimensions so that it can make a picture that refers to the byte map (in this
case 4) which it can scale, rotate, whatever. The backend gets told that it has
to include that byte map instead of the picture (path) so again \LUA\ kicks in
and we get \in {figure} [fig:integration:2].

\startplacefigure[reference=fig:integration:2]
    \scale[width=.5tw]{\getbuffer}
\stopplacefigure

Especially in the \LUAMETAFUN\ manual there are various examples of the tight
integration of the three languages, with occasional help from a \CCODE\ library.
We wrap up with two examples where we use a small library that generates
a \QRCODE.

\startbuffer[1]
\startluacode
    context.dontleavehmode()
    figures.qrcode(
        "To use LUA or METAPOST, that's the question.",
        tex.sp("5cm")
    )
\stopluacode
\stopbuffer

\startbuffer[2]
\startluacode
local colors = {
    string.char(200,  0,  0),
    string.char(  0,200,  0),
    string.char(  0,  0,200),
    string.char(  0,200,200),
    string.char(200,  0,200),
    string.char(200,200,  0),
}

local white = string.char(255,255,255)
local black = string.char(  0)

function MP.Test(index,str)
     local data, size = qrcodegen.generate(str)
     data = string.gsub(data,"(.)",function(c)
        if c == black then
            return colors[math.random(6)]
        else
            return white
        end
     end)
     mp.newbytemap(index,size,size,3,data)
end
\stopluacode
\stopbuffer

\startbuffer[3]
\startMPcode
    lua.MP.Test(3, "To use LUA or METAPOST, that's the question.") ;
    draw bytemap 3 xsized 5cm ;
\stopMPcode
\stopbuffer

The first variant used a \LUA\ helper in the \CONTEXT\ graphic library. Here
\METAPOST\ is not involved:

\typebuffer[1][option=TEX]

The second example calls a \LUA\ function from the \METAPOST\ end where we
replace the black rectangles with colorful ones. My phone can handle it.

\typebuffer[2][option=TEX]
\typebuffer[3][option=TEX]

The reason why we added this library is that you can configure \CONTEXT\ to pop
up an error message in \QRCODE\ speak. We also used these codes on the \CONTEXT\
2025 meeting poster and for ear tags with familiar \TEX\ error messages at
\BACHOTEX. Anyway, we show both examples side by side in \in {figure}
[fig:integration:3]. I hope that we made clear that the possibilities of
combining techniques are endless.

\startplacefigure[reference=fig:integration:3]
    \startcombination
        {\getbuffer[1]}   {Just \LUA.}
        {\getbuffer[2,3]} {\LUA\ and \METAPOST.}
    \stopcombination
\stopplacefigure

\stopchapter

\stopcomponent

% This time's musical timestamp: Leprous Live in Tilburg Bluray (2025)
